\section{Problem Formulation and Superstitious-Memory Hypothesis}

\paragraph{Setting and notation.}
An agent processes a stream of tasks $\{(q_t, r_t)\}$, where $q_t$ is a query or
task and $r_t \in [0,1]$ is an outcome signal (an environment success indicator or
a consistency-checked judge score). It maintains a memory bank
$\bank_t = \{m_i\}$, a time/state-indexed retrieval operator
$\ret_t(q_t;\bank_t,\phibar,S) = C_t \subseteq \bank_t$
returning a top-$k$ set (the operator depends on the current bank, running credit,
and scope state, so it changes over deployment), and a policy $\policy(y \given q, C)$ that produces an
output $y$. \method{} adds an attributor $\attr$, a per-memory running credit
$\phibar_i$, and a scope representation $S_i$. These are the core symbols;
Algorithm~\ref{alg:merit} lists the full runtime state (label pools, trajectories,
per-memory counts, prototypes, quarantine and calibration state).

\paragraph{Immediate retrieval contribution.}
We define a memory's value by intervention, at a single retrieval event and
conditioned on the deployment history.

\begin{definition}[Immediate retrieval contribution]\label{def:phi}
For a retrieval event with query $q$, realized retrieved set $C=\ret(q)$, and
history $H_t$,
\begin{equation}
\phicf_{i,t}(q,C) \triangleq
  \E_\xi\!\left[\, r_t(C;\xi) - r_t(C\setminus\{m_i\};\xi) \given H_t \,\right],
\end{equation}
the expected outcome difference from removing $m_i$, over policy and environment
randomness $\xi$.
\end{definition}

\noindent This is the single-event effect, \emph{not} a long-horizon memory value;
long-horizon value is out of scope. The estimand is the \emph{literal removal}
effect. Because \rit{} implements the removal by inserting a length-matched neutral
pad rather than leaving an empty slot (Section~4), the interventional labels may be
interpreted as estimating the removal effect \emph{only if} the neutral-pad
equivalence validation (\tbd{PILOT_NEUTRAL_PAD_VALIDATION}) passes a pre-registered
equivalence test; if that validation fails or has not been run, the estimand
downgrades to the \emph{pad-replacement contribution}, and literal removal remains
unvalidated future work. Higher-order interactions within $C$ are not part of this
estimand: leave-one-out is the target, and where interactions matter they are
surfaced only in the boundary study through an exploratory group-intervention
analysis (A7) that carries no numerical guarantee (AS3).
We separate four
related quantities and never conflate them: the true effect $\phicf$; its
finite-rollout \rit{} estimate $\phirit$ (Definition~\ref{def:phi} evaluated with
$K$ paired rollouts), which estimates the removal effect under the equivalence
precondition above; the amortized \aca{} prediction $\phihat$; and the online
running aggregate $\phibar_i$ (an EMA of $\phihat$). At the memory level we
distinguish the true aggregate target $\phiagg_i \triangleq \E[\phicf_{i,t}]$ from
its finite \rit{} audit estimate $\phiaggrit_i$ (reported with a confidence
interval; Section~4). The aggregation is taken over an explicit \emph{audit
distribution}: the retrieval-conditioned event distribution of $m_i$ on frozen bank
snapshots, restricted to the common-support region, with uniform weight per audited
retrieval event; $\phiaggrit_i$ is the corresponding design-based (probability-sampled)
finite estimator. The three memory-level signals $\coutil_i$, $\phibar_i$, and
$\phiaggrit_i$ are compared on this same audit population. The window, minimum sample,
and inclusion design are locked before running (\emph{USER\_APPROVAL\_REQUIRED};
Section~4). The baseline signal, by contrast, is associational.

\begin{definition}[Co-occurrence utility]\label{def:util}
The baseline memory-level signal is
\begin{equation}
\coutil(m_i) = \frac{\sum_t \mathbf{1}[m_i \in \ret(q_t)]\, r_t}
                    {\sum_t \mathbf{1}[m_i \in \ret(q_t)]},
\end{equation}
the average outcome of tasks in which $m_i$ was retrieved.
\end{definition}

\paragraph{Why correlational credit can be superstitious.}
\begin{observation}[Decomposition identity]\label{obs:decomp}
Because $\phicf$ is a definitional removal difference, the memory-level
decomposition
\begin{equation}
\E[r(C)\given i\text{ retr.}] =
  \underbrace{\E[r(C\setminus\{m_i\})\given i\text{ retr.}]}_{\text{leave-one-out baseline}}
  + \underbrace{\E[\phicf_i\given i\text{ retr.}]}_{\text{mean contribution}}
\end{equation}
holds as an \emph{identity}, requiring no assumption. Consequently $\coutil_i$ mixes
the leave-one-out baseline success with the mean contribution: it is a
\emph{composite signal that can be misaligned with immediate contribution},
differing from $\E[\phicf_i\given i\text{ retrieved}]$ by exactly the leave-one-out
baseline term, which is nonzero whenever
$\E[r(C\setminus\{m_i\})\given i\text{ retr.}]\neq 0$.
\end{observation}

\noindent The identity is not conjectural. It does not by itself establish that a
\emph{persistent} misalignment arises in a running system; that is a separate,
genuinely conjectural claim about the closed-loop dynamics, which we state as a
hypothesis rather than a theorem.

\begin{conjecture}[Superstition equilibrium, informal]\label{conj:equilibrium}
If retrieval is statistically dependent on baseline solvability, this dependence
amplifies and systematizes the misalignment; and under positive feedback in which
utility participates in retrieval ranking, a stable \emph{superstition equilibrium}
can arise, with $\coutil_i$ high while $\phicf_i = 0$.
\end{conjecture}

\noindent Retrieval--solvability dependence is a mechanism that may amplify the
misalignment under the conditions tested in the pilot; it is \emph{not} necessary
for the misalignment to exist or for the decomposition identity
(Observation~\ref{obs:decomp}) to hold. A formal statement and proof of the
equilibrium claim are deferred; Conjecture~\ref{conj:equilibrium} is not proven in
this paper and must not be read as established.

\paragraph{Assumptions.}
Three assumptions scope the method and are stated for rebuttal. To avoid collision
with the ablation labels A1--A7 used later, we name the assumptions AS1--AS3.
\begin{itemize}
  \item \textbf{AS1 (observable outcome).} Each task admits a usable $r_t$, either
  an environment success signal or a consistency-checked judge score.
  \item \textbf{AS2 (local replayability).} A task subset can be replayed from a
  fixed initial state (e.g., container snapshots or text environments). Only
  \rit{} label collection requires AS2; \aca{} online scoring does not.
  \item \textbf{AS3 (interaction caveat).} Leave-one-out is the estimand; it does
  not capture higher-order interactions within a retrieved set, and we do not treat
  such interactions as estimation error against any allocation rule (e.g., the
  Shapley value). Where interactions matter---most visibly for
  redundant/substitutable memory pairs, whose individual leave-one-out effects can
  be near zero while their joint value is positive---we surface them only in the
  boundary study through an exploratory group-intervention analysis (A7) that
  carries no numerical guarantee.
\end{itemize}

\paragraph{Objectives.}
\aca{} learns $\attr$ by minimizing a regression loss on \rit{} labels,
$\E[(\attr(x_i(q)) - \phirit_i(q))^2]$, where $x_i(q)$ is a feature vector for the
event. The system objective maximizes streaming cumulative success $\sum_t r_t$
subject to a token budget: retrieval plus \rit{} sampling stay within
$1 + \varepsilon$ times the baseline cost, with
$\varepsilon = \tbd{RIT_BUDGET_EPSILON_PERCENT}$ as a provisional hyperparameter.
